\begin{prob}

        It turns out that BFS does not need to use three distinct statuses
        (undiscovered, pending, visited) in order to
        function -- we just use three to help us understand the algorithm.  In
        fact, we can delete code from bfs, and the
        function will still work. List below the line number(s) which can all
        be deleted without affecting the behavior of the algorithm.

        \begin{soln}
            Line 26\\
            Note that we need \textbf{some} intermediate step (i.e, the 'pending' state) so that the check on
            line 22 evaluates to False for the nodes that we've discovered but have not visited yet. 
            Otherwise, we would always add the neighbors of \texttt{u} to the \texttt{pending} queue, even if they
            have already been discovered. As a result, the outer while loop will undergo unnecessary iterations on each call to BFS.\\

            However, we never explictly check for a node being marked as 'visited', and the code would execute 
            the same way if we left all discovered nodes as 'pending'.
        \end{soln}

\end{prob}
